% \begin{table*}[!ht]
% \centering
% \caption{Behavioural gap results for ASTER and EvoSuite across four projects. Beh. Cov.\% is the proportion of expected behaviours validated by the tool's suite. Beh. Gap\% is the conservative lower-bound estimate of confirmed behavioural gaps using the developer-written confirmation rate of 70.6\%. \rth{can we get this down to a single column? we care about project, tool, methods, UTBs and behavioural gaps. if we drop behaviours and beh. coverage, is it narrow enough? if not, we could just make it two tables and drop the tool column, or build the tool name into a row in the table instead of being a full column.}}
% \label{tab:rq3}
% \begin{tabular}{llrrrrr}
% \toprule
% \textbf{Project} & \textbf{Tool} & \textbf{Methods} & \textbf{Behaviours} & \textbf{UTBs} & \textbf{Beh. Cov.\%} & \textbf{Beh. Gap\%} \\
% \midrule
% commons-cli      & ASTER    & 169 & 333   & 91  & 72.7\% & 19.3\% \\
% commons-codec    & ASTER    & 275 & 656   & 255 & 61.1\% & 27.4\% \\
% commons-compress & ASTER    & 200 & 371   & 165 & 55.5\% & 31.4\% \\
% commons-jxpath   & ASTER    & 134 & 271   & 108 & 60.1\% & 28.1\% \\ \midrule
% \textbf{Total}   & & \textbf{778} & \textbf{1,631} & \textbf{619} & \textbf{62.0\%} & \textbf{26.8\%} \\
% \midrule
% \midrule
% commons-cli      & EvoSuite & 213 & 420   & 122 & 71.0\% & 20.5\% \\
% commons-codec    & EvoSuite & 470 & 1,091 & 318 & 70.9\% & 20.6\% \\
% commons-compress & EvoSuite & 812 & 1,406 & 440 & 68.7\% & 22.1\% \\
% commons-jxpath   & EvoSuite & 292 & 543   & 120 & 77.9\% & 15.6\% \\ \midrule
% \textbf{Total}   &  & \textbf{1,787} & \textbf{3,460} & \textbf{1,000} & \textbf{71.1\%} & \textbf{20.4\%} \\
% \bottomrule
% \end{tabular}
% \end{table*}



\subsection{RQ3 (Generalization): To what extent are auto\-matically-generated test suites prone to behavioural gaps?}
\label{sec:rq3}

This research question seeks to investigate whether behavioural gaps are inherent to developer-written suites, or a systemic limitation of structure-guided testing practices.
To do this, we evaluate whether behavioural gaps persist in test suites created by automatic test generation approaches.
We chose two state‑of‑the‑art (SOTA) automated test generators: EvoSuite~\cite{10.1145/2025113.2025179} and ASTER~\cite{Pan2024ASTERNA}, to evaluate whether they can systematically close behavioural gaps.
ASTER was selected because it represents the SOTA for LLM-supported test generation, and is a shipping industrial product for automatically generating test suites.
EvoSuite was selected because it is the baseline test generation tool widely used in academia, and was also used by the ASTER team in their own comparative evaluation.
Evaluating whether automatically generated test suites would help contextualize whether behavioural gaps are just a byproduct of misplaced human effort, or a more pervasive testing challenge. 

\smallskip
\noindent\textbf{Methodology:} 
ASTER provides a robust replication package including both ASTER and EvoSuite generated tests for eight projects.
Unfortunately, ASTER is not publicly available so we restricted our analysis to the four projects in the ASTER dataset that intersected with the projects in our dataset.
To enable a fair comparison between the three tools, we only evaluated ASTER-\toolname{} on public documented methods ASTER was able to generate tests for. The same holds for the methods examined in the EvoSuite-\toolname{} comparison. 
This accounts for the different method counts in Table~\ref{tab:rq3}.

To ensure behavioural gaps were not just due to test suites having low structural metric values, we also computed the line coverage for the analyzed methods used for each test generation tool.
For ASTER, 81.2\% of methods achieve 91–100\% line coverage; for EvoSuite, 89.7\% achieve 91-100\% line coverage.

Since the four projects were already in our dataset, we already had the expected behaviours for all of the methods being analyzed, but \toolname{} had to be executed again with the test suites generated by ASTER and EvoSuite to identify the behavioural gaps in these suites.
For this analysis we only report the estimated number of behavioural gaps by multiplying the number of UTBs identified by \toolname{} by the 0.713 precision from RQ2.


\begin{table}[h]
\centering
\caption{UTB and method concentration across coverage buckets (aggregated across 
all projects). UTBs (\%) is the share of total UTBs falling in each bucket. 
Methods (\%) is the share of total methods in each bucket.}
\label{tab:rq4}
\small
\setlength{\tabcolsep}{3.5pt}
\begin{tabular}{l cc cc}
\toprule
& \multicolumn{2}{c}{\textbf{Line Coverage}} 
& \multicolumn{2}{c}{\textbf{Mutation Score}} \\
\cmidrule(lr){2-3} \cmidrule(lr){4-5}
\textbf{Bucket} 
& \textbf{UTBs (\%)} & \textbf{Methods (\%)} 
& \textbf{UTBs (\%)} & \textbf{Methods (\%)} \\
\midrule
% \multicolumn{5}{l}{\textit{Developer-written tests}} \\
$\leq$70\%  & 5.2  & 2.8  & 44.3 & 24.3 \\
71--80\%    & 4.1  & 2.4  & 3.0  & 3.1  \\
81--90\%    & 5.4  & 2.6  & 2.6  & 2.7  \\
91--100\%   & 85.4 & 92.3 & 50.1 & 69.8 \\
% \midrule
% % ppp: Not sure if we should keep the ASTER and Evosuite here as those are only for the 4 projects, for now commenting it out
% \multicolumn{5}{l}{\textit{ASTER-generated tests}} \\
% $\leq$70\%  & 21.5 & 11.8 & 56.5 & 40.2 \\
% 71--80\%    & 4.7  & 3.9  & 0.7  & 2.3  \\
% 81--90\%    & 4.2  & 3.1  & 1.0  & 0.8  \\
% 91--100\%   & 69.6 & 81.2 & 41.8 & 56.7 \\
% \midrule
% \multicolumn{5}{l}{\textit{EvoSuite-generated tests}} \\
% $\leq$70\%  & 23.1 & 7.6  & 74.7 & 69.2 \\
% 71--80\%    & 1.8  & 1.4  & 0.6  & 0.6  \\
% 81--90\%    & 1.2  & 1.3  & 0.5  & 0.5  \\
% 91--100\%   & 73.9 & 89.7 & 24.2 & 29.8 \\
\bottomrule
\end{tabular}
\end{table}


\smallskip
\noindent\textbf{Estimating UTB prevalence in generated suites:}
%
The behavioural gaps left by both ASTER and EvoSuite are reported in Table~\ref{tab:rq3}.
In summary, 27.1\% of expected behaviours were not validated by the ASTER-generated suites and 20.6\% of the expected behaviours were not validated by the EvoSuite-generated suites.
These levels of behavioural gaps both exceed the developer-written suites for these four projects, as detailed in Table~\ref{tab:rq2_dev}.

This suggests that neither EvoSuite, a search-based test generator, nor ASTER, an LLM-based test generator are immune to missing expected behaviours in software systems.
One reason for this could be because neither ASTER nor EvoSuite use documentation to guide their test generation. 

We also examined Judot~\cite{denaro2025automatedtestgenerationprogram}, a documentation-aware test generator that uses JDoctor to extract pre- and postconditions from structured Javadoc tags. 
Across the four projects, it generated tests for only 149 methods, compared to 778 for ASTER and 1,787 for EvoSuite, since its parser is limited to formal tags and ignores narrative documentation. 
On the methods it did reach, the tests showed a 42.5\% behavioural gap, suggesting that even documentation-aware generation, when restricted to structured tags, captures only a fraction of the behaviours developers express in method-level documentation.


%To ameliorate this, we investigated a documentation-aware test generator approach called Judot~\cite{denaro2025automatedtestgenerationprogram} which extracts formal pre- and postconditions from structured Javadoc tags to drive test generation.  
%However, this design imposes a fundamental reach limitation imposed by Judot's underlying JDoctor implementation. JDoctor can only extract conditions from structured tags, leaving behaviours expressed in narrative text entirely out of scope. 

%As noted in Section~\ref{sec:related}, this limitation is inherent to structured documentation parsers. In practice, Judot exercised only 149 methods across the four projects, compared to 778 for ASTER and 1,787 for EvoSuite.
%This was because JDoctor was unable to extract conditions for the remaining methods. 
%Even among the 149 methods it does reach, \toolname identified 192 UTBs out of 319 expected behaviours, yielding a conservative behavioural gap estimate of 42.5\%. This suggests that structured contract extraction, while useful for what it captures, covers only a fraction of what developers actually document.

% JDoctor was unable to generate conditions for most of the identified methods and hence Judot couldn't generate any tests for them. Among the 149 methods it does reach, \toolname identified 192 UTBs out of identified 319 expected behaviour. Applying the conservative lower bound yields an estimated behavioural gap rate of 42.5\%. Even on the narrow set of methods where Judot succeeds, more than half of the documented behaviours remain unvalidated. This suggests structured contract extraction while useful for what it captures covers only a fraction of what developers actually document.

\begin{table}[!t]
\centering
\caption{Behavioural gap results for ASTER and EvoSuite across four projects. 
%Beh. Cov.\% is the proportion of expected behaviours validated by the tool's suite. 
Beh. Gap\% is the conservative lower-bound estimate of confirmed behavioural gaps using the developer-written confirmation rate of 71.3\%.}
\label{tab:rq3}
\begin{tabular}{llrrr}
\toprule
\textbf{Project} & \textbf{Tool} & \textbf{Methods} & \textbf{UTBs} & \textbf{Beh. Gap\%} \\
\midrule
commons-cli      & ASTER    & 169   & 91  & 19.5\% \\
commons-codec    & ASTER    & 275 & 255 & 27.7\% \\
commons-compress & ASTER    & 200 & 165 & 31.7\% \\
commons-jxpath   & ASTER    & 134 & 108 & 28.4\% \\ \midrule
\textbf{Total}   & & \textbf{778} & \textbf{619} & \textbf{27.1\%} \\
\midrule
\midrule
commons-cli      & EvoSuite & 213  & 122 & 20.7\% \\
commons-codec    & EvoSuite & 470 & 318 & 20.8\% \\
commons-compress & EvoSuite & 812 & 440 & 22.3\% \\
commons-jxpath   & EvoSuite & 292 & 120 & 15.8\% \\ \midrule
\textbf{Total}   &  & \textbf{1,787} & \textbf{1,000} & \textbf{20.6\%} \\
\bottomrule
\end{tabular}
\vspace{-2em}
\end{table}
\smallskip
\noindent\textbf{Implication:} Both of the test generation approaches failed to have smaller behavioural gaps than the developer-written suites. 
This suggests that behavioural gaps do not arise due to developer inattention, but instead because they capture aspects of software validation that are orthogonal to the structural metrics test generation approaches traditionally optimize for.

\begin{rqbox}{RQ3: Prevalence of generated behavioural gaps.}
Behavioural gaps are not unique to developer‑written suites. Suites generated by ASTER and EvoSuite exhibit behavioural gaps of 27.1\% and 20.6\% respectively, despite achieving strong structural coverage on the methods they exercise.
\end{rqbox}